\documentclass[12pt]{extarticle}
\def\activityid{F02-U-v3}\usepackage[letterpaper,margin=.65in,top=.60in,bottom=.65in]{geometry}
\usepackage{fix-cm}
\usepackage[T1]{fontenc}\usepackage[default]{sourcesanspro}
\usepackage{microtype,tikz,array,tabularx,fancyhdr,amsmath,amssymb}
\usepackage[hidelinks]{hyperref}\usetikzlibrary{calc,arrows.meta}
\setlength{\parindent}{0pt}\setlength{\parskip}{7pt}\setlength{\headheight}{0pt}\setlength{\footskip}{26pt}
\setlength{\emergencystretch}{2em}\pagestyle{fancy}\fancyhf{}\renewcommand{\headrulewidth}{0pt}\renewcommand{\footrulewidth}{.4pt}
\fancyfoot[L]{\scriptsize Bellingham Math Circle / Week 2 / \activityid}\fancyfoot[R]{\scriptsize \thepage}
\definecolor{ink}{gray}{.12}\definecolor{soft}{gray}{.95}\color{ink}
\newcommand{\PageTitle}[2]{{\fontsize{10}{12}\selectfont\bfseries WEEK 2\quad TOGGLE LAB\hfill #1\par}\vspace{3pt}{\fontsize{26}{29}\selectfont\bfseries #2\par}\vspace{2pt}}
\newcommand{\NameLine}{{\small Name: \rule{2.65in}{.4pt}\hfill Date: \rule{1.35in}{.4pt}\par}\vspace{4pt}}
\newcommand{\Task}[2]{\par\vspace{3pt}{\large\bfseries #1}\par #2\par}
\newcommand{\Subhead}[1]{\par\vspace{3pt}{\large\bfseries #1}\par}
\newcommand{\RuleBox}[1]{\par\begingroup\setlength{\fboxsep}{8pt}\colorbox{soft}{\parbox{\dimexpr\linewidth-16pt\relax}{#1}}\endgroup\par}
\newcommand{\WriteBox}[2]{\par\textbf{#1}\par\vspace{2pt}\begin{tikzpicture}\draw[black!40,line width=.5pt] (0,0) rectangle (\dimexpr\linewidth-.5pt\relax,#2);\end{tikzpicture}\par}
\newcommand{\StudentFooter}{\vfill{\small Try it with objects. Explain by pointing, talking, or drawing.}\par}
\newcommand{\LampRing}[3]{\begin{tikzpicture}[x=1in,y=1in]
\foreach \i in {1,...,#1}{\coordinate (v\i) at ({90-360*(\i-1)/#1}:#2);}
\foreach \i in {1,...,#1}{\pgfmathtruncatemacro{\j}{mod(\i,#1)+1}\draw[thick] (v\i)--node[midway,fill=white,inner sep=2pt,font=\small]{\i\j}(v\j);}
\foreach \i in {1,...,#1}{\node[circle,draw,fill=white,minimum size=.52in,inner sep=0pt,font=\bfseries] at(v\i){\i};}
\foreach \i in {#3}{\node[circle,draw,fill=ink,text=white,minimum size=.52in,inner sep=0pt,font=\bfseries] at(v\i){\i};}
\end{tikzpicture}}
\newcommand{\Lights}[1]{\begin{tikzpicture}[x=.55in,y=.55in,baseline=-.10in]\foreach \v [count=\i] in {#1}{\ifnum\v=1\fill (\i,0) circle(.16in);\else\draw (\i,0) circle(.16in);\fi}\end{tikzpicture}}
\newcommand{\ClockRing}[2]{\begin{tikzpicture}[x=1in,y=1in]\draw[black!30] (0,0) circle(#2);\pgfmathtruncatemacro{\n}{#1-1}\foreach\i in {0,...,\n}{\node[circle,draw,fill=white,minimum size=.35in,inner sep=0pt] at({90-360*\i/#1}:#2){\i};}\draw[-{Stealth},thick] (-.35,.15) arc(155:25:.38);\node[font=\small] at(0,-.18){clockwise};\end{tikzpicture}}
\newcommand{\Key}[2]{\begin{center}\renewcommand{\arraystretch}{1.55}\begin{tabular}{c|*{#1}{c}}in&#2\end{tabular}\end{center}}


% v3 student pages use one running header and numbered tasks only.
\newcommand{\StudentHeader}[1]{{\fontsize{10}{12}\selectfont\bfseries Week 2 / Toggle lab / #1\par}\vspace{10pt}}
\newcommand{\Problem}[2]{\par\vspace{4pt}\textbf{Problem #1:} #2\par}
\newcommand{\Space}[1]{\begin{tikzpicture}\draw[black!35] (0,0) rectangle (\dimexpr\linewidth-.5pt\relax,#1);\end{tikzpicture}\par}
\newcommand{\Record}[2]{#1\par\Space{#2}}

\begin{document}
\StudentHeader{Grades 4--5}

\Problem{1}{Start with five two-sided counters OFF. Pressing a line flips both end lamps, including ON to OFF. Write 12 for pressing the line from 1 to 2. Try two moves and undo them.}
\begin{center}\LampRing{5}{1.06}{}\end{center}
\Problem{2}{Start all OFF for each target. Make only the listed lamps ON. Find two different move lists, then circle your shorter list.}
\begin{center}\renewcommand{\arraystretch}{2.4}\begin{tabular}{l|p{2.05in}|p{2.05in}}ON lamps&First list&Another list\\\hline 1 and 2&&\\1 and 3&&\\1, 2, 3, 4&&\\\end{tabular}\end{center}
\Record{Choose a target. Test whether one press can make it. Record what happened.}{.75in}

\newpage
\StudentHeader{Grades 4--5}

\Problem{3}{Start all OFF for each list. Try 12,23; then 23,12; then 12,23,12,12. Draw the final lamps for each. Cross out any pair of repeated presses that you can remove without changing the final picture.}
\Space{1in}
\Problem{4}{A reduced list uses each edge at most once. For each reduced list below, make its target. Reset; press exactly the edges missing from that list. Compare the two final pictures and count the presses. The five edges are 12,23,34,45,51.}
\begin{center}\renewcommand{\arraystretch}{2.5}\begin{tabular}{l|p{1.65in}|p{1.6in}|p{.8in}}First list&Missing edges&Final ON lamps&Counts\\\hline 12&&&\\12,23&&&\\12,34&&&\\12,23,34&&&\\\end{tabular}\end{center}
\Problem{5}{Choose a new edge set and its missing edges. Test both. Record the two press counts and whether the final targets agree.}
\Space{1.05in}

\newpage
\StudentHeader{Grades 4--5}

\Problem{6}{Use the five-lamp mat. To leave only lamps 1 and 3 ON, first choose whether to press 51. Lamp 1 then tells you whether to press 12. Keep going around the ring. Fill yes or no, test the full list, and check lamp 5 last.}
\begin{center}\renewcommand{\arraystretch}{2.2}\begin{tabular}{c|*{5}{p{.47in}}|p{.95in}}First choice&51&12&23&34&45&Lamp 5\\\hline No&no&&&&&\\Yes&yes&&&&&\\\end{tabular}\end{center}
\Problem{7}{Repeat the forced choices for only lamps 1,2,3,4 ON. Compare your two lists with your experiments on page 1.}
\begin{center}\renewcommand{\arraystretch}{2.2}\begin{tabular}{c|*{5}{p{.47in}}|p{.95in}}First choice&51&12&23&34&45&Lamp 5\\\hline No&no&&&&&\\Yes&yes&&&&&\\\end{tabular}\end{center}
\Problem{8}{Now try to leave only lamp 1 ON. Fill both rows. Put a cross beside any row whose final lamp does not match the target.}
\begin{center}\renewcommand{\arraystretch}{2.2}\begin{tabular}{c|*{5}{p{.47in}}|p{.95in}}First choice&51&12&23&34&45&Lamp 5\\\hline No&no&&&&&\\Yes&yes&&&&&\\\end{tabular}\end{center}
\Record{Record a rule you notice about the two lists or their targets.}{.9in}

\newpage
\StudentHeader{Grades 4--5}

\Problem{9}{On the five-lamp ring, choose three new targets with an even number of ON lamps. Find the shortest reduced list for each. Record both complementary lists and their lengths. Try to find a target needing more than two presses.}
\Space{1.25in}
\Problem{10}{Use six counters on this ring. Start all OFF. Make only lamps 1 and 4 ON in two different ways. Count the presses. Then repeat for only 1 and 3 ON. Compare the two lengths for each target.}
\begin{center}\LampRing{6}{1.06}{}\end{center}
\Record{Move lists, lengths, and a comparison with the five-lamp ring:}{1.0in}

\end{document}
